% =====================================================================
% 考研政治 · 重点总结手册 —— 正文
% =====================================================================

\zzchapter{第一部分\quad 马克思主义基本原理}

\zzsection{一、基本知识}

\zzsub{科学社会主义的核心命题}

\zzcard{科学社会主义的核心命题}{%
  \kw{“是2必”}：\bk{资本主义必然灭亡，社会主义必然胜利}。\\
  理论依据：资本主义基本矛盾（生产社会化与生产资料资本主义私人占有之间的矛盾）决定资本主义必然被社会主义取代。
}

\zzsub{世界的物质统一性原理}

\zzcard{世界的物质统一性原理}{%
  世界的物质统一性原理是\\
  \bk{辩证唯物主义最基本、最核心的观点}\\
  \bk{马克思主义的基石}
}

\zzsub{马恩第一个无产阶级政党}

\zzcard{马恩第一个无产阶级政党}{%
  \bk{马克思和恩格斯}建立的第一个\bk{无产阶级政党}：\kw{共产主义者同盟}。
}

\zzsub{一个错误说法}

\zzcard{一个错误说法}{%
  理论联系实际，\bk{不是}“实践符合理论”。\\
  \zzwarn{用\bk{实践套理论}＝\kw{本本主义}；用\bk{实践发展理论}＝\kw{中国化}。}
}

\zzsub{理解政治经济学的枢纽}

\zzcard{理解政治经济学的枢纽}{%
  \kw{“牢二枢纽”}：理解政治经济学的枢纽是\bk{劳动二重性理论}。\\
  \zzwarn{辨析：马克思主义经济学说的\bk{核心内容和基石}是\kw{剩余价值理论}。}
}

\zzsub{六个寓于}

\zzcard{六个寓于}{%
  \begin{zztable2}
    命题 & 助记 \\
    \midrule
    价值寓于使用价值之中 & （J>S）驾驭 \\
    必然性寓于偶然性之中 & （B>O）碧玉 \\
    斗争性寓于同一性之中 & （D>T）斗鱼 \\
    绝对性寓于相对性之中 & （J>X）绝育 \\
    普遍性寓于特殊性之中 & （P>T）璞玉 \\
    共性寓于个性之中 & （G=G）公寓 \\
  \end{zztable2}
  \zzwarn{“寓于”\bk{前面}的一定是\bk{绝对}的东西，\bk{无条件}的要寓于\bk{有条件}的。}
}

\zzsub{马克思主义的基本特征}

\zzcard{马克思主义的基本特征}{%
  \bk{本质属性}：\kw{人民性}\\
  \bk{根本特征（显著特征）}：\kw{实践性}\\
  \bk{首要的、基本的观点}：\kw{实践的观点}\\
  \zzwarn{一句话概括：\kw{科学性与革命性的统一}（科命）。}
}

\zzsub{唯物史观和剩余价值学说}

\zzcard{唯物史观和剩余价值学说}{%
  唯物史观和剩余价值学说是\\
  \bk{社会主义由空想变为科学的两大基石}\\
  \bk{马克思一生的两个伟大发现}
}

\zzsub{人的发展}

\zzcard{人的发展}{%
  人的发展是一个社会历史过程，既包括人类整体的发展，也包括人类个体的发展。\\
  人的发展指人的体力、智力、个性和交往能力等的发展，其中\bk{最根本的是人的自由程度的提高}。\\
  \zzwarn{无论社会或个人，\bk{根本目标}都是：\kw{人的自由全面的发展}。}
}

\zzsub{对立统一规律}

\zzcard{对立统一规律}{%
  对立统一规律是\kw{唯物辩证法的实质和核心}，这是因为：\\
  \begin{enumerate}
    \item 它揭示了事物普遍联系的\bk{根本内容}和变化发展的\bk{内在动力}，从根本上回答了事物为什么会发展的问题；
    \item 它是贯穿\bk{质量互变规律}、\bk{否定之否定规律}以及唯物辩证法基本范畴的\bk{中心线索}，也是理解这些规律和范畴的\bk{钥匙}；
    \item 它提供了人们认识世界和改造世界的根本方法——\kw{矛盾分析方法}。
  \end{enumerate}
}

\zzsub{马克思主义的三大组成部分}

\zzcard{马克思主义的三大组成部分}{%
  \bk{马克思主义哲学}是\kw{基础和方法}\\
  \bk{马克思主义政治经济学}是\kw{主体内容}\\
  \bk{科学社会主义}是\kw{目的和归宿}\\[2pt]
  《德意志意识形态》标志着\bk{马哲}的诞生\\
  《资本论》标志着\bk{政经}的诞生\\
  《共产党宣言》标志着\bk{科社}的诞生和整个马克思主义的诞生\\
  《反杜林论》首次将马克思主义以\bk{三个组成部分}的形式进行了全面阐述\\
  \zzwarn{马克思主义科学的世界观和方法论基础是\kw{辩证唯物主义和历史唯物主义}。}
}

\zzsub{马克思主义相关著作}

\zzcard{马克思主义相关著作}{%
  \begin{zztable2}
    著作 & 核心结论 \\
    \midrule
    《德法年鉴》 & 完成从唯心主义向唯物主义、从革命民主主义向共产主义的\bk{转变}，为创立马克思主义奠定思想前提（德法两转变） \\
    《神圣家族》 & \bk{马克思恩格斯}合著的第一部著作，批判黑格尔维新，论述历史唯物（马恩第一家族） \\
    《德意志意识形态》 & 首次系统阐述了\bk{历史唯物主义}的基本观点（德历史） \\
    《共产党宣言》 & 发表标志着马克思主义的\bk{公开问世}（公共） \\
    《资本论》 & 被称为\bk{“工人阶级的圣经”}（工资） \\
    《法兰西内战》 & 科学总结了\bk{巴黎公社}的历史经验（法国巴黎） \\
    《哥达纲领批判》 & 对未来社会从制度、发展阶段以及分配等方面作了原则构想（哥够香） \\
    《反杜林论》 & 全面阐述了马克思主义理论体系，被称为马克思主义的\bk{“百科全书”}（百度） \\
  \end{zztable2}
  \zzsub{恩格斯《1848年至1850年的法兰西阶级斗争》导言}
  恩格斯去世前为该书写导言，指出：1.\ 社会主义革命的胜利将是一个\bk{长期的过程}；2.\ 利用资产阶级民主和普选权的\bk{重要性}；3.\ 在重视合法斗争的同时，\bk{没有否定暴力革命}的作用，恩格斯晚年没有放弃革命。\\[2pt]
  \zzwarn{（伯恩施坦）修正主义：反对无产阶级革命，主张通过和平的\bk{议会斗争}和改良措施，使资本主义“和平长入”社会主义。}
}

\zzsection{二、资本}

\zzsub{金融资本形成的主要途径}

\zzcard{金融资本形成的主要途径}{%
  \kw{“金铲铲”}：\bk{金融联系}、\bk{资本参与}、\bk{人事参与}。
}

\zzcard{货币最基本的职能}{%
  \kw{是“甲流”}：\bk{价值尺度和流通手段}。
}

\zzcard{货币职能}{%
  \bk{支付手段}——债务、赋税、工资、租金。
}

\zzsub{资本主义国家内部的矛盾}

\zzcard{资本主义国家内部的矛盾}{%
  \begin{enumerate}
    \item \bk{经济发展失调}：①虚拟经济与实体经济发展失衡；②福利风险增加；③债务负担沉重。
    \item \bk{政治体制失灵}：①选举；②政党利益；③民主陷阱；④精英政治。
    \item \bk{社会融合机制失效}：①社会极端思潮抬头；②社会流动性退化；③社会矛盾激化。
  \end{enumerate}
}

\zzsub{金融寡头控制形式}

\zzcard{金融寡头控制形式}{%
  \begin{enumerate}
    \item 金融寡头在\bk{经济}领域实现统治的形式是\kw{“参与制”}（股票）；
    \item 金融寡头在\bk{政治}领域中实现对\bk{国家机器}的控制主要是通过同政府的\kw{“个人联合”}；
    \item 金融寡头左右和影响内政外交与社会生活，主要是通过\bk{建立政策咨询机构}，掌握新闻、科教、文化等。
  \end{enumerate}
}

\zzsub{国际金融垄断资本主义}

\zzcard{国际金融垄断资本主义}{%
  \kw{“空心瓜果”}：\bk{国际金融资本的垄断}成为当代资本主义\bk{最突出、最鲜明、最主要的特征}，表现在：\\
  \bk{金融垄断寡头化}、\bk{金融垄断国际化}、\bk{经济虚拟化、产业空心化}。\\
  \zzwarn{金融垄断资本形成壮大的重要制度条件：\kw{金融自由化与金融创新}。在金融垄断资本的推动下，\bk{虚拟经济越来越脱离实体经济}。}
}

\zzsub{垄断资本向世界范围扩展的基本形式}

\zzcard{垄断资本向世界范围扩展的基本形式}{%
  \kw{“垄扩借产品”}：\bk{借贷资本输出}、\bk{生产资本输出}、\bk{商品资本输出}。\\
  \zzwarn{辨析——\bk{资本在循环过程中的三种职能}：货币资本、生产资本、商品资本。}
}

\zzsub{所有制辨析与资本周转}

\zzcard{所有制辨析}{%
  \bk{私人股份所有制}：\kw{所有权和控制权不再统一}。\\
  \bk{法人资本所有制}：\kw{“狼狠独揽大权”}。
}

\zzcard{资本周转快慢}{%
  影响资本周转快慢的关键因素有两个：\\
  ①资本周转的\bk{时间}；②生产资本的\bk{固定资本和流动资本的构成}。\kw{（构原）}
}

\zzsub{资本积累}

\zzcard{资本积累}{%
  \bk{资本积累的唯一源泉}就是\kw{剩余价值}；\\
  资本积累是\bk{资本主义扩大再生产的源泉}；\\
  资本积累导致了\bk{资本主义社会的两极分化}；\\
  资本积累的扩大\bk{会导致资本有机构成提高}，进而\bk{产生失业问题}；\\
  资本积累是\bk{资本主义制度走向灭亡的自我否定因素}。\\
  \zzwarn{资本有机构成的提高是一般趋势，由资本的本性决定。\bk{资本有机构成一直在增加，看到说减少的就是错的}。}
}

\zzsub{生产资本的划分方式}

\zzcard{生产资本的划分方式}{%
  \begin{zztable2}
    划分依据 & 类别 \\
    \midrule
    根据资本是否增值（是否带来剩余价值） & 不变资本、可变资本 \\
    根据在资本循环中的不同职能 & 货币资本、生产资本、商品资本 \\
    根据资本周转方式不同 & 固定资本、流动资本 \\
  \end{zztable2}
}

\zzsub{不变资本和可变资本}

\zzlongcard{不变资本和可变资本}{%
  根据\bk{资本在剩余价值生产中所起的不同作用}划分：\\
  ①【\bk{不变资本}】是以\bk{生产资料}形态存在的资本；\\
  ②【\bk{可变资本}】是用来\bk{购买劳动力}的那部分资本。\\[2pt]
  \bk{可变资本的价值}是\kw{工人的劳动再生产出来的}（可变资本是工人干活干出来的）；\\
  \bk{不变资本的价值}是\kw{具体劳动转移到新产品中}（可一次转移，也可多次转移）。\\
  \zzwarn{“……的价值”前面，\bk{物的价值被转移，人的价值被创造}。}\par\addvspace{2pt}
  【\bk{剩余价值}】是雇佣工人所创造并被资本家无偿占有的超过劳动力价值的那部分价值，是雇佣工人剩余劳动的凝结。\\
  技术不断提高、资本积累增长、资本积聚和集中，个别资本规模日益增长：\bk{机器变多，生产资料（不变资本）变多，人变少}，\bk{有机构成（机器人）逐渐变多，“可变资本”相对占比逐渐变少}。\\
  \zzwarn{技术变，价值变，有机也就变；技术不变，价值变，有机就不变。\bk{失业无法消除}。}
}

\zzsub{资本主义工资}

\zzcard{资本主义工资}{%
  资本主义工资表现为\bk{劳动的价值或价格}，这是一种\kw{假象}。\\
  实际上，工人出卖的不是劳动——\bk{劳动不是商品}，工人出卖的是\bk{劳动力}，\bk{劳动力是商品}。\\
  因此，资本主义工资\bk{本质上}是\kw{劳动力的价值或价格}，而不是劳动的价值或价格。
}

\zzsub{矛盾}

\zzcard{三类矛盾}{%
  ①\bk{资本主义的基本矛盾}：\kw{生产社会化和生产资料资本主义私人占有之间的矛盾}。\\
  ②\bk{私有制基础上商品经济的基本矛盾}：\kw{私人劳动和社会劳动的矛盾}。它：√决定着\bk{商品经济的本质及发展过程}；√是\bk{商品经济其他一切矛盾的基础}；√决定着\bk{商品生产者的命运}。\\
  ③商品的使用价值和价值的矛盾、生产商品的具体劳动和抽象劳动的矛盾，\bk{根源于私人劳动和社会劳动的矛盾}。\\
  ④\bk{社会形态更替}归根结底是社会基本矛盾运动的结果，\bk{生产力的发展具有最终的决定意义}。
}

\zzsection{三、生产和价值}

\zzsub{社会再生产的核心问题}

\zzcard{社会再生产的核心问题}{%
  \bk{社会再生产的核心问题是社会总产品的实现问题}，即社会总产品的\kw{价值补偿和实物补偿}问题。\kw{（生孩子之前要补偿她）}
}

\zzsub{社会}

\zzcard{社会}{%
  在社会生产中，【\bk{生产力}】是生产的【\bk{物质内容}】，【\bk{生产关系}】是生产的【\bk{社会形式}】，二者构成【\bk{社会的生产方式}】。\\
  \bk{社会形态更替}归根结底是【\bk{社会基本矛盾运动}】的结果，其中\bk{生产力的发展具有最终的决定意义}。
}

\zzsub{价值规律的作用}

\zzcard{价值规律的作用}{%
  价值规律在市场配置资源过程中的作用表现在：\\
  （1）自发地调节\bk{生产资料和劳动力}在社会各生产部门之间的\bk{分配比例}；\\
  （2）自发地刺激\bk{社会生产力的发展}；\\
  （3）自发地调节\bk{社会收入的分配}。\\
  \zzwarn{× 自发地调节雇佣工人与资本家之间的利益再分配——\bk{错误}。资本家压榨工人（干活多、发工资少）这种现象\bk{不会被价值规律所调节}。}
}

\zzsub{超额和相对剩余价值}

\zzcard{超额和相对剩余价值}{%
  单个资本家改进技术、改善管理的\bk{主观动机}是追求\kw{超额剩余价值}，但其\bk{客观后果}则是整个社会各个生产部门的劳动生产率普遍提高，导致生活资料的价值下降和补偿劳动力价值的必要劳动时间缩短，剩余劳动时间相对延长，\bk{整个资本家阶级普遍获得更多的相对剩余价值}。\\
  \zzwarn{助记：\kw{超额对应原因、个别}；\kw{相对对应结果、社会}。}
}

\zzsub{劳动生产率}

\zzcard{劳动生产率}{%
  \bk{劳动生产率}指的是劳动者生产\bk{使用价值}的效率。它的高低可以用\bk{单位劳动时间内生产的商品数量}来测量，也可以用\bk{单位商品中所耗费的劳动时间}来测量。\\
  通俗理解：①单位时间内你做得\bk{多不多}；②单位商品你\bk{多久做出来}（快不快）。\\[2pt]
  \bk{商品的价值}是凝结在商品中的无差别的一般人类劳动，价值量由\bk{劳动者生产商品所耗费的劳动量}决定，而劳动量按\bk{劳动时间}计量。因此，\kw{在单位劳动时间内，劳动者耗费的劳动量不变，生产的价值总量也不变}。
}

\zzsub{垄断价格}

\zzcard{垄断价格}{%
  \bk{垄断高价}是指\kw{卖的贵}；\bk{垄断低价}是指\kw{买的便宜}，即进原材料便宜、正常价卖。\\
  \zzwarn{垄断高价和垄断低价都是\bk{赚钱}的价格，不会低于原材料价格，也抵消不了。}\par\addvspace{2pt}
  垄断价格产生之后，\bk{从全社会看商品的价格总额等于商品的价值总额}。\\
  \bk{垄断利润的形成，关键在于垄断组织在经济生活中起了决定性作用}。\\
  第一，垄断没有消除产生竞争的经济条件；第二，垄断必须通过竞争来维持；第三，社会生产是复杂多样的，任何垄断组织都不可能把包罗万象的社会生产全部包下来。\\
  在竞争手段上，自由竞争主要运用\bk{经济手段}，垄断条件下的竞争除了经济手段外，还采取\bk{非经济的手段}。
}

\zzsection{四、政经计算}

\zzsub{核心公式}

\zzcard{政经计算}{%
  \begin{zztable2}
    符号 / 公式 & 含义 \\
    \midrule
    不变资本 $C$、可变资本 $V$、剩余价值 $M$ & 三大基本量 \\
    投资＝成本＝$C+V$ & 预付总资本 \\
    有机构成 & $C:V$ \\
    剩余价值率 & $M:V$ \\
    利润 & $=M$（利润与剩余价值在量上一样） \\
  \end{zztable2}
  \zzsub{典型算例}
  某资本家投资 100 万元，即 $c+v=100$ 万元；其预付资本的\bk{有机构成是 $3:1$}，即 $c:v=3:1$，则 $c=75$ 万元，$v=25$ 万元。\\
  该资本家每次投资的\bk{剩余价值率为 200\%}，即 $m/v=200\%$，则每次投资的剩余价值 $m=50$ 万元。\\
  利润与剩余价值在量上一样，故每次投资的\bk{利润为 50 万元}。
}

\zzsub{绝对和相对剩余价值计算}

\zzcard{绝对和相对剩余价值计算}{%
  √\bk{绝对剩余价值}生产方法，比较\bk{大胆}，总的工作时间\bk{延长}；\\
  √\bk{相对剩余价值}生产方法，比较\bk{隐蔽}，工作日时间\bk{不变}。\\
  \zzwarn{绝对剩余价值的生产和相对剩余价值的生产都是依靠\bk{延长工人的剩余劳动时间}获得的。}\par\addvspace{2pt}
  \bk{剩余价值率＝剩余劳动时间÷必要劳动时间}。\\
  若原工作时间 8 小时，必要劳动时间 4 小时，剩余劳动时间 4 小时：\\
  \begin{zztable2}
    方法 & 计算 \\
    \midrule
    绝对剩余价值（延长劳动时间） & 时间 $8+2=10$，剩余价值率 $=6\div4\times100\%=\textbf{150\%}$ \\
    相对剩余价值（缩短必要劳动时间） & 必要劳动时间 $4-2=2$，剩余价值率 $=6\div2\times100\%=\textbf{300\%}$ \\
  \end{zztable2}
}

\zzsection{五、国际（列宁、战时共产主义与新经济政策、巴黎公社）}

\zzsub{列宁}

\zzlongcard{列宁}{%
  ① 列宁指出，\bk{资本主义发展到帝国主义，在经济生活中起决定作用的是垄断组织}；\\
  ② 列宁提出：\bk{社会主义革命可以首先在一个或者几个国家获得胜利}；\\
  ③ 在俄国社会主义革命取得胜利的初期，特别是实行新经济政策期间，列宁提出了许多精辟论述：\\
  \hspace{1em}1.\ 把建设社会主义作为一个\bk{长期探索、不断实践}的过程；\\
  \hspace{1em}2.\ 把\bk{大力发展生产力、提高劳动生产率}放在首要地位；\\
  \hspace{1em}3.\ 在多种经济成分并存的条件下，\bk{利用商品、货币和市场}发展经济；\\
  \hspace{1em}4.\ \bk{利用资本主义建设社会主义}。\\
  列宁还指出了\bk{马克思主义执政党建设的极端重要性}，阐明了\bk{思想建设和文化建设}的重要意义，提出了加强国家政权建设和发扬社会主义民主的一系列措施。\\
  ④ 列宁在被称为“政治遗嘱”的文章和书信中提出建设社会主义的新构想，主要包括：\\
  \hspace{1em}用\bk{合作社}的形式将农民引向社会主义道路；发展大工业，实现\bk{工业化和电气化}；学习和利用\bk{资本主义}一切有价值的东西；进行\bk{文化革命}，大力发展文化教育事业；进行党和国家机构的改革，努力提高干部的素质和能力；必须\bk{反对官僚主义}，健全社会主义民主和法制；维护党的团结，特别是党中央领导核心的团结。
}

\zzsub{战时共产主义与新经济政策}

\zzcard{三种模式对比}{%
  \begin{zztable3}
    模式 & 核心内容 & 目的 \\
    \midrule
    战时共产主义 & 余粮收集制、取消商品货币关系 & 更好地供给军队作战 \\
    新经济政策（列宁） & 恢复商品货币关系，发展资本主义 & 战后迅速活跃经济 \\
    苏联模式（斯大林） & 农业合作化 & 进一步发展社会主义 \\
  \end{zztable3}
  \zzwarn{列宁为社会主义\bk{奠定了基础}，但列宁\bk{并未发展}社会主义。}\par\addvspace{2pt}
  \zzsub{新经济政策（与战时共产主义比较）}
  ①\bk{农业}：\kw{粮食税制}取代余粮收集制——不再全部收走，交完税后多的归自己，允许农民自由出卖余粮；\\
  ②\bk{工业}：部分恢复\bk{私营经济}，允许私人经营中小企业，但关系国家命脉的企业仍归国家所有；实行\bk{国家资本主义}，采取\bk{租让制和租借制}，允许外资企业管理国家暂时无力经营的企业；\\
  ③\bk{贸易}：\kw{允许自由贸易}（废除了国家配给和国家贸易垄断制）；\\
  ④\bk{分配}：实行\kw{按劳分配}；\\
  ⑤\bk{流通}：\kw{恢复商品货币关系}。
}

\zzsub{巴黎公社}

\zzcard{巴黎公社}{%
  1871 年爆发的巴黎公社革命，是\bk{第一国际精神的产儿}，是\kw{无产阶级夺取政权的第一次伟大尝试}。\\
  特别是采取了\bk{两项重要措施}，用以防止国家机关工作人员由社会公仆变为社会主人：\\
  一是规定所有公职人员无论职位高低，实行全面的\bk{选举制和撤换制}；二是\bk{取消高薪制}，规定任何工作人员年薪不得超过熟练工人的工薪水平。\\[2pt]
  \bk{要点}：①巴黎公社建立的是\bk{政权}，不是\bk{政党}；②防止公仆变主人的两项措施\bk{与人相关}，即人的工资和职位——取消高薪、选举撤换。\\
  \zzwarn{辨析：\bk{俄国十月革命}才建立\bk{苏维埃大会制度}。巴黎公社革命源于民族矛盾和阶级斗争激化，\bk{不是第一国际发动和领导}，只是\bk{受第一国际精神影响}。}
}
